\documentclass[11pt,a4paper]{article}
% Estilo comum dos PDFs individuais da Sprint 1.
\usepackage[utf8]{inputenc}
\usepackage[T1]{fontenc}
\usepackage[brazilian]{babel}
\usepackage[a4paper,margin=2.35cm]{geometry}
\usepackage{xcolor}
\usepackage{hyperref}
\usepackage{enumitem}
\usepackage{listings}
\usepackage{parskip}
\usepackage{needspace}
\definecolor{azul}{RGB}{20,63,112}
\hypersetup{colorlinks=true,linkcolor=azul,urlcolor=azul}
\setlist[itemize]{leftmargin=*,itemsep=3pt}
\lstset{basicstyle=\ttfamily\small,frame=single,breaklines=true,columns=fullflexible}
\sloppy
\newcommand{\codigo}[1]{\texttt{#1}}
\newcommand{\titulo}[2]{\begin{center}{\LARGE\bfseries #1\par}\vspace{4pt}{\large #2\par}\vspace{8pt}\rule{0.8\textwidth}{0.4pt}\par\vspace{5pt}{\small Matheus Moreli Durão Silva -- Matrícula 12621CBS217\par Cibersegurança - Bacharelado - Noturno\par}\end{center}}
\newcommand{\cardfonte}[1]{\par\smallskip\noindent\textit{Referência: }\href{#1}{Trello - enunciado e checklist do card}.}

\begin{document}\titulo{S1\_13 - Enumerações}{Estados válidos para um ativo}
\section*{Entregável correto}
O card pede o exercício de status de ativos com \codigo{Enum}. Uma enumeração evita que cada parte do programa escreva uma variação diferente de um mesmo estado, como ``ativo'', ``Ativo'' ou ``em uso'' sem padrão.
\section*{Como criar e acessar uma enumeração}
Importa-se \codigo{Enum}, declara-se uma classe e cada membro recebe valor. O membro é acessado pelo nome, por exemplo \codigo{StatusAtivo.ATIVO}; o valor associado pode ser mostrado com \codigo{.value}. Ao comparar, compare membros da mesma enumeração, não textos improvisados.
\begin{lstlisting}
from enum import Enum

class StatusAtivo(Enum):
    ATIVO = "Ativo"
    MANUTENCAO = "Em manutencao"
    INATIVO = "Inativo"

status = StatusAtivo.ATIVO
if status == StatusAtivo.ATIVO:
    print(status.value)
for opcao in StatusAtivo:
    print(opcao.value)
\end{lstlisting}
\section*{Aplicação no projeto}
O \codigo{app.py} usa \codigo{IntEnum} para categorias de ativos, mantendo opções numeradas consistentes. Para salvar um \codigo{Enum} em JSON, não se grava o objeto diretamente; salva-se \codigo{.value} ou um valor numérico e, ao ler, reconstrói-se a enumeração. Isso explica por que padronização e persistência precisam ser pensadas juntas.
\section*{Checklist demonstrado}
Conceito de Enum, criação, acesso a membros e valores, comparação, iteração e contexto do projeto. O teste é percorrer todos os status e conferir que o ativo só recebe uma das opções definidas.
\cardfonte{https://trello.com/c/nJrQKljh}
\end{document}
